\section{Root Extraction}
\subsection*{SQUARE ROOT 1 --- B4-10.0}
\label{sec:square_root_1}

\subsubsection*{PURPOSE}
The program \texttt{B4-10.0} calculates the square root of a positive number. The number must reside in the accumulator at an even $q$ factor. The program utilizes Newton's method to solve the equation:
\begin{equation}
    0 = x^2 - N
\end{equation}
by means of successive approximations:
\begin{equation}
    x_{i+1} = x_i + \frac{1}{2} \left( \frac{N}{x_i} - x_i \right)
\end{equation}
The program operates as a subroutine and is invoked via a calling sequence in the main program. Upon branching to the subroutine, the argument must reside in the accumulator, and the return address must be stored within the subroutine. Upon return to the main program, the calculated square root is held in the accumulator.

\subsubsection*{STORAGE REQUIREMENTS}
\textbf{Program Space:} 51 words \\
\textbf{Temporary Storage:} None

\subsubsection*{INPUT}
The input variable for ``Square Root 1'' is $N$, the positive number whose square root is to be calculated.

$N$ must reside in the accumulator with an even $q$ factor when branching to the subroutine.

\subsubsection*{CALLING SEQUENCE}
\begin{table}[h!]
\centering
\begin{tabularx}{\textwidth}{cccX}
\toprule
\textbf{Memory Location} & \textbf{Instruction} & \textbf{Address} & \textbf{Meaning} \\
\midrule
$\alpha-1$ & \texttt{B} & \texttt{L} & Bring $N$ into the accumulator at an even $q$. \\
$\alpha$ & \texttt{R} & \texttt{L+50} & Store the return address in the subroutine. \\
$\alpha+1$ & \texttt{U} & \texttt{L} & Unconditional branch to the subroutine entry. \\
\bottomrule
\end{tabularx}
\end{table}

This function uses a standard LGP-21 function call. The only notable factor is that $N$ must have an even $q$.\footnote{See \label{subsub:halfQ} for why.}

\subsubsection*{OUTPUT}
Upon the return branch to the main program, the square root of the argument $N$ resides in the accumulator at $q_N / 2$. For example, if $N$ is at $@24$, then $\sqrt{N}$ will be at $@12$.

\subsubsection*{OPERATING DETAILS}

\begin{enumerate}[leftmargin=1.5em, itemsep=6pt]
    \item \textbf{Accuracy:} \\
    No error is greater than $2^{-30}$.

    \item \textbf{Execution Time:} \\
    Approximately 775~ms.

    \item \textbf{Program input:} \\
    The program tape contains the routine in two formats:
    \begin{itemize}
      \item Relocatable hexadecimal, valid for J1-10.0 (section \ref{sec:PIRloader})
        \item Relocatable decimal, in coding sheet format.
    \end{itemize}

    \item \textbf{Required Input/Output Units:} \\
    None.

    \item \textbf{Overflow:} \\
    Overflow is ignored and is not reset upon return.

    \item \textbf{Error Stops:} \\
    \begin{tabularx}{\textwidth}{lX}
        \toprule
        \textbf{Memory Location} & \textbf{Meaning and Remedy} \\
        \midrule
        \texttt{L+24}            & $N$ is negative. After pressing the \textsf{START} button, the accumulator is cleared, and control returns directly to the main program. \\
        \bottomrule
    \end{tabularx}
\end{enumerate}
\newpage
\subsection*{$n$-TH ROOT --- B4-11.0}
\label{sec:nth_root}

\subsubsection*{PURPOSE}
The program \texttt{B4-11.0} calculates the $n$-th root of any non-negative value. The exponent $n$ must be an integer such that $2 \le n \le 20$. The results are calculated using an iterative method according to the formula:
\begin{equation}
    y_{i+1} = y_i + \frac{1}{n} \left( \frac{A}{y_i^{n-1}} - y_i \right)
\end{equation}
where $y_i$ represents the $i$-th approximation of the root.

\noindent
\texttt{B4-11.0} uses a calling sequence similar to \ref{sec:logarithm_function}, except it has only one secondary parameter. The number $A$ whose root is to be taken is loaded into the accumulator with an appropriate scale factor (see below), and the degree of the root $n$ is stored in a \texttt{Z} instruction following the \texttt{R} that calls the subroutine. On return to the main program, the desired root resides in the accumulator at $@q/n$.

\subsubsection*{STORAGE REQUIREMENTS}
\textbf{Program Space:} 1 track (64 words) \\
\textbf{Temporary Storage:} Sectors 01, 02, 04, 05, 44, 53, 54, and 60 of track 63.

\subsubsection*{INPUT}
The input variables for \texttt{B4-11.0} are $A$, the non-negative value whose root is to be calculated, and $n$, the integer specifying the degree of the root ($2 \le n \le 20$). 

Upon branching to the subroutine, $A$ must reside in the accumulator scaled to a $q$ factor that is exactly divisible by $n$.\footnote{This is similar to the requirement that a number whose square root is being taken is divisible by two; however, $n$ only needs to divide $A$'s $q$ factor exactly to be acceptable. As an example, if we had $A$ at $@12$, then the square, cube, fourth, and sixth roots of $A$ could all be taken, as all are even divisors of 12.}

\subsubsection*{CALLING SEQUENCE}
\begin{table}[h!]
\centering
\begin{tabularx}{\textwidth}{cccX}
\toprule
\textbf{Memory Location} & \textbf{Instruction} & \textbf{Address} & \textbf{Meaning} \\
\midrule
$\alpha - 1$ & \texttt{B} & $A$ & Bring $A$ into accumulator at a $q$ divisible by $n$. \\
$\alpha$ & \texttt{R} & \texttt{L+16} & Store the parameter address in the subroutine. \\
$\alpha + 1$ & \texttt{U} & \texttt{L} & Unconditional branch to the subroutine entry. \\
$\alpha + 2$ & \texttt{Z} & \texttt{08nn} & Parameter $n$ encoded as a decimal sector address. \\
\bottomrule
\end{tabularx}
\end{table}

The instruction at location $\alpha - 1$ brings the value $A$ into the accumulator at a $q$ factor that is divisible by $n$. (Any instruction sequence that achieves this result is permissible.)

The instruction at location $\alpha$ stores the address ($\alpha + 2$) containing the parameter instruction for $n$ within the subroutine, and execution resumes at $\alpha+3$.\footnote{The subroutine takes the parameter address supplied via the \texttt{R} instruction, recalculates the return address, and adjusts its own code to properly return.}

\subsubsection*{OUTPUT}
Upon the return branch to the main program, the $n$-th root of the argument resides in the accumulator scaled to a $q$ factor of $q_A / n$. For example, if $n = 3$ and $q_A = 12$, the result $\sqrt[3]{A}$ will have a $q$ of at $@4$.

\subsubsection*{OPERATING DETAILS}

\begin{enumerate}[leftmargin=1.5em, itemsep=6pt]
    \item \textbf{Accuracy:} \\
    The result is guaranteed accurate to eight decimal places.
    
    \item \textbf{Execution time:} \\
    Between 2 and 60 seconds. Execution times are significantly longer when calculating high-degree roots of small values.
    
    \item \textbf{Required input/output devices:} \\
    None.
    
    \item \textbf{Overflow:} \\
    Overflow is ignored upon entry and is not reset upon exit.
    
    \item \textbf{Error stops:} \\
    \begin{tabularx}{\textwidth}{lX}
        \toprule
        \textbf{Memory Location} & \textbf{Meaning and Remedy} \\
        \midrule
        \texttt{L+61}            & The requested root is an imaginary value (i.e., an even-degree root of a negative number). Pressing the \textsf{START} button clears the accumulator and returns execution directly to the main program. \\
        \bottomrule
    \end{tabularx}
\end{enumerate}
\newpage